\section{Оптимизация программы}

Для устранения выявленного «узкого горлышка» архитектура дерева была полностью переписана:
\begin{enumerate}
    \item \textbf{Переход на итеративный подход:} рекурсивные вызовы \texttt{Insert}, \texttt{Remove} и \texttt{Search} были заменены на циклы \texttt{while}. Для выполнения балансировки «снизу вверх» путь от корня до текущего узла сохраняется в локальный массив \texttt{std::vector<TNode**>}. Резервирование памяти (\texttt{reserve(64)}) исключает динамические аллокации во время спуска.
    \item \textbf{Отказ от умных указателей:} \texttt{std::unique\_ptr} были заменены на «сырые» (raw) указатели \texttt{TNode*}. Управление памятью теперь осуществляется вручную (удаление дерева при вызове деструктора словаря).
\end{enumerate}

\section{Сравнение результатов (До и После)}

Программа была перекомпилирована с включённым профилировщиком. Выдержка из отчёта \texttt{gprof} для оптимизированной версии:
\begin{alltt}
  %   cumulative   self              self     total           
 time   seconds   seconds    calls  ns/call  ns/call  name    
 16.67      0.21     0.05 21951385     2.28     2.28  std::max (UpdateHeight)
  6.67      0.26     0.02 19951425     1.00     8.11  std::vector::push_back()
  6.67      0.28     0.02 18951425     1.06     1.06  std::operator> (string)
  ...
  0.00      0.30     0.00  1000000     0.00   270.00  AVLTree::Insert()
\end{alltt}
Вызовы умных указателей полностью исчезли из профиля. Теперь основное процессорное время затрачивается на полезную работу: пересчёт высот (\texttt{std::max}) и сравнение строк. Среднее время одной вставки упало с \textbf{910 нс} до \textbf{270 нс}.

Для чистого замера производительности обе версии были скомпилированы с флагом \texttt{-O3} без профилировщиков. Использовалась системная утилита \texttt{time} (метрика \texttt{user time}). 

\begin{table}[h!]
\centering
\begin{tabular}{|l|c|c|l|}
\hline
\textbf{Тестовый набор (1 млн оп.)} & \textbf{До (Рекурсия)} & \textbf{После (Итерация)} & \textbf{Ускорение} \\ \hline
Random Insertion & 1.207 с & \textbf{1.080 с} & $\approx 1.1x$ \\ \hline
Sorted Insertion & 0.295 с & \textbf{0.260 с} & $\approx 1.1x$ \\ \hline
Search           & 0.336 с & \textbf{0.271 с} & $\approx 1.2x$ \\ \hline
\end{tabular}
\caption{Сравнение «чистого» времени выполнения релизных версий}
\label{tab:bench_results}
\end{table}

Пиковое потребление памяти (Massif) после перехода на сырые указатели незначительно снизилось и составило \textbf{68.58 MB}. Отсутствие рекурсии также позволило избежать риска переполнения стека (\texttt{Stack Overflow}) при крайне несбалансированных деревьях во время интенсивных удалений.
\pagebreak